\section{Introduction}
\label{sec:intro}

Longitudinal data, in which the same entities are followed over time, is now routine across the sciences and across public policy. Countries are tracked year by year through their emissions and their income, monitoring stations record air quality month by month, and household surveys follow families through their budgets. A common question in all of these settings is not only how a quantity moves on average, but whether the population hides distinct paths, and if so, who follows each path and why. Answering that question well calls for two things at once. The distinct paths have to be recovered from the data, and each path has to be described in terms a domain reader can act on.

Existing tools tend to supply one of these and not the other. Trajectory clustering recovers distinct paths, yet it hands back cluster membership without a rule that would let a reader recognise the members of a cluster from their attributes. Subgroup discovery supplies exactly such rules, yet its usual target is a single number attached to each entity rather than a shape traced over a shared calendar. A reader who wants both, a recovered path and a rule that names the entities on it, is left to bolt the two together by hand, which is fragile and hard to reproduce.

We introduce SEAL, a method for interpretable trajectory segmentation that produces both at once. Each entity contributes a target trajectory over a shared time grid. SEAL groups the trajectories by how they accumulate over time and describes every group by a short conjunctive rule over the entity attributes. It gives the two halves equal weight and searches for the grouping and its rule together, so the description is part of what the search optimises rather than a report added afterwards. A candidate grouping survives only when a simple rule captures it, and a rule is kept only when it tracks a coherent trajectory shape. Because each description is a shallow decision rule, the membership of any group reads as a single sentence.

A concrete policy question motivates the work. When an unexpected illness strikes, a family can be forced to spend a large share of its income on health, and knowing which households see that share rise or fall, and on the back of what earlier event, matters for social protection. Household budget panels also test an interpretable method sharply, because the natural target is a ratio of health spending to income, and a careless analysis rediscovers the definition of that ratio and reports it as a finding. We treat the household health panel as the flagship study and accompany it with four public panels, covering carbon dioxide emissions, pandemic mortality, national income, and air quality, whose driving structure is documented and can therefore check the discovered rules.

We make three contributions. First, we introduce SEAL, which couples two tasks that are usually kept apart, grouping trajectories by shape and describing each group by a rule, and joins them through a single cost that neither half can satisfy alone. Second, we turn covariates that vary over time into engineered temporal descriptors, so that a rule can state how a covariate moved over the window rather than only where it ended, and we show when this sharpens the discovered groups, most clearly where the structure turns on a covariate's movement. Third, we impose a non-circularity discipline on the descriptors, so that a rule cannot restate the definition of its own target, which is essential when the target is a ratio such as a health-spending share. Across all five panels the discovered rules select groups that match documented structure, and on the household panel they identify who carries the burden of health spending and reveal how that burden moved through the pandemic.

The rest of the paper is organised as follows. Section~\ref{sec:related} reviews the lines of work that we draw on. Section~\ref{sec:problem} states the problem. Section~\ref{sec:algorithm} describes the method and the descriptor engineering. Section~\ref{sec:data} sets out the data and the configuration. Section~\ref{sec:cases} presents the five case studies, and Section~\ref{sec:conclusion} concludes.
